\section*{Introduction}
\addcontentsline{toc}{section}{Introduction}

Ce document est une annexe au rapport de projet de la fusée bi-étage SP-02, développée par
l'association AéroIPSA au cours de l'année scolaire 2025-2026 \cite{rapportSP02}. Il n'a pas
vocation à présenter le projet dans son ensemble : l'architecture générale de la fusée, le
dimensionnement de ses sous-systèmes et le déroulement de la mission y sont supposés connus.
La lecture préalable du rapport principal est donc nécessaire à la bonne compréhension du
présent document, qui n'en développe qu'un point particulier.\\

Pour situer brièvement le sujet traité ici, SP-02 est une fusée expérimentale à deux étages
dont la séparation intervient en vol, après l'extinction du propulseur du premier étage. Le
mécanisme qui assure cette séparation doit satisfaire deux exigences antagonistes : maintenir
une liaison rigide entre les deux corps pendant toute la phase propulsée, où il transmet
l'intégralité des efforts d'inertie et de flexion subis par la fusée, puis libérer cette
liaison de manière franche sur ordre du séquenceur. Il constitue à ce titre un élément
critique de la mission : sa défaillance en vol compromettrait aussi bien la tenue structurale
de l'ensemble qu'une séparation propre des étages.\\

Ce mécanisme repose sur deux pièces principales, une pièce de liaison en aluminium 7075-T6
reliant le tube du premier étage au système de connexion du second, et un jeu de loquets
assurant le verrouillage mécanique de l'interface. Ce sont ces pièces, et elles seules, qui
font l'objet de la présente étude.\\

Aucune campagne d'essais mécaniques instrumentés ne pouvant être menée dans le cadre du
projet, la simulation numérique constitue le principal moyen de vérifier le dimensionnement
de ces pièces avant leur usinage. L'analyse par \acrfull{FEM} a donc été conduite en parallèle
de leur conception, afin d'évaluer leur tenue sous les charges de vol et d'identifier
d'éventuelles zones critiques.\\

L'approche \acrshort{FEM} présente pour ce projet plusieurs intérêts. Elle donne accès aux
champs de contraintes et de déformations dans des géométries que les modèles analytiques
classiques ne savent pas décrire, en raison des congés, perçages et variations d'épaisseur
qui les caractérisent. Elle permet de localiser les concentrations de contraintes autour des
fixations et des poches d'allègement, zones où s'amorcent en pratique les défaillances. Elle
autorise enfin l'étude du comportement dynamique des pièces par le calcul de leurs fréquences
propres. Ces résultats sont systématiquement confrontés à des calculs analytiques simplifiés,
qui servent de garde-fou sur les ordres de grandeur.\\

Le document s'organise en deux parties. La première décrit le cadre méthodologique retenu :
la chaîne logicielle Patran-Nastran, les étapes de préparation et de résolution des modèles,
le protocole de calcul analytique servant de référence, les paramètres d'étude et les limites
inhérentes à la méthode. La seconde présente les analyses menées sur le mécanisme de
séparation, pièce par pièce, puis confronte les résultats numériques aux estimations
analytiques.

\newpage